% ══════════════════════════════════════════════
% CHƯƠNG 4 — THỰC THI KỸ THUẬT VÀ CÔNG NGHỆ
% Người viết: TV4 (API — Controllers, JWT, Swagger)
% ══════════════════════════════════════════════

\chapter{Thực thi Kỹ thuật và Công nghệ}

\section{ASP.NET Core Web API}\label{Sec4.1}

\subsection{Thiết kế RESTful Endpoints}\label{sub4.1.1}

\begin{de}\label{def4.1}
  % [TV4] Định nghĩa REST: các HTTP verb chuẩn (GET/POST/PUT/PATCH/DELETE),
  % nguyên tắc đặt tên URL (danh từ, không dùng động từ), status codes.

  % [Điền nội dung vào đây]
\end{de}

% [TV4] Điền toàn bộ endpoint thực tế của hệ thống vào bảng dưới.
% Thêm/bớt dòng tùy theo số endpoint thực tế.

\begin{longtable}{|l|l|l|l|}
  \hline
  \textbf{Method} & \textbf{Endpoint} & \textbf{Mô tả} & \textbf{Quyền} \\
  \hline
  % GET    & /api/[resource]           & [Mô tả]          & [Public/Role] \\
  % GET    & /api/[resource]/\{id\}    & [Mô tả]          & [Public/Role] \\
  % POST   & /api/[resource]           & [Mô tả]          & [Public/Role] \\
  % PUT    & /api/[resource]/\{id\}    & [Mô tả]          & [Public/Role] \\
  % DELETE & /api/[resource]/\{id\}    & [Mô tả]          & [Public/Role] \\
  % ... thêm các endpoint khác
  \hline
  \caption{Danh sách các RESTful Endpoints chính}\label{table:endpoints}
\end{longtable}

\subsection{Cấu trúc Controller}\label{sub4.1.2}

% [TV4] Chèn code C# của 1 Controller tiêu biểu, bao gồm ít nhất
% GET (list), GET (detail), POST, PUT, DELETE.

\begin{lstlisting}[style=csharp, caption={[TênController] — các action chính}]
// [TV4] Dán code Controller vào đây
\end{lstlisting}

\section{Xác thực và Phân quyền (JWT + RBAC)}\label{Sec4.2}

\subsection{JWT Authentication}\label{sub4.2.1}

\begin{thm}\label{thm4.1}
  % [TV4] Phát biểu về JWT: stateless, thông tin được mã hóa trong token,
  % không cần lưu session phía server.

  % [Điền nội dung vào đây]
\end{thm}

\noindent\begin{proof}
  % [TV4] Giải thích tại sao JWT đảm bảo stateless và dễ scale ngang
  % (server chỉ verify chữ ký, không tra DB).

  % [Điền nội dung vào đây]
  \QEDAa
\end{proof}

% [TV4] Chèn code cấu hình JWT trong Program.cs
\begin{lstlisting}[style=csharp, caption={Cấu hình JWT trong Program.cs}]
// [TV4] Dán code AddAuthentication / AddJwtBearer vào đây
\end{lstlisting}

\subsection{Phân quyền theo Role}\label{sub4.2.2}

% [TV4] Mô tả các Role trong hệ thống và quyền hạn tương ứng.
% Điền vào bảng phân quyền dưới đây.

\begin{table}[H]
  \begin{center}
    \begin{tabular}{|l|c|c|c|}
      \hline
      \textbf{Chức năng} & \textbf{[Role 1]} & \textbf{[Role 2]} & \textbf{[Role 3]} \\
      \hline
      % [Chức năng 1] & \checkmark & \checkmark &            \\
      % [Chức năng 2] & \checkmark &            &            \\
      % [Chức năng 3] & \checkmark & \checkmark & \checkmark \\
      % ... thêm dòng nếu cần
      \hline
    \end{tabular}
  \end{center}
  \caption{Bảng phân quyền theo Role}\label{table:phanquyen}
\end{table}

\section{Global Exception Middleware}\label{Sec4.3}

% [TV4] Mô tả ngắn mục đích của Middleware xử lý lỗi tập trung,
% sau đó chèn code.

% [Điền nội dung vào đây]

\begin{lstlisting}[style=csharp, caption={ExceptionMiddleware xử lý lỗi tập trung}]
// [TV4] Dán code ExceptionMiddleware vào đây
\end{lstlisting}

\section{LINQ to XML — Export và Import}\label{Sec4.4}

\begin{ex}\label{ex4.1}
  % [TV4] Mô tả tính năng Export/Import XML bằng LINQ to XML:
  % dùng ở đâu trong hệ thống, dữ liệu gì được export/import?

  % [Điền nội dung vào đây]
\end{ex}

% [TV4] Chèn code Export XML
\begin{lstlisting}[style=csharp, caption={Export dữ liệu ra XML dùng LINQ to XML}]
// [TV4] Dán code tạo XDocument / XElement vào đây
\end{lstlisting}

% [TV4] Chèn code Import XML
\begin{lstlisting}[style=csharp, caption={Import dữ liệu từ XML dùng LINQ to XML}]
// [TV4] Dán code parse XDocument vào đây
\end{lstlisting}

\section{Kiểm thử API với Swagger và Postman}\label{Sec4.5}

% [TV4] Chụp màn hình Swagger UI đang chạy, lưu thành file
%       rồi thay \framebox bên dưới bằng: \includegraphics[width=14cm]{swagger.png}

\begin{figure}[H]
  \centering
  \framebox[14cm][c]{\rule{0pt}{6cm} [Chèn screenshot Swagger UI]}
  \caption{Giao diện Swagger UI kiểm thử API}\label{fig:swagger}
\end{figure}

% [TV4] Chụp màn hình Postman kết quả test (Collection Runner hoặc request/response),
%       lưu thành file rồi thay \framebox bên dưới.

\begin{figure}[H]
  \centering
  \framebox[14cm][c]{\rule{0pt}{5cm} [Chèn screenshot Postman test results]}
  \caption{Kết quả kiểm thử API bằng Postman}\label{fig:postman}
\end{figure}

\section{Dashboard (nếu có)}\label{Sec4.6}

% [TV4/TV5] Nếu nhóm xây dựng giao diện Frontend/Dashboard,
% mô tả ngắn công nghệ sử dụng và chèn screenshot.
% Nếu không có, có thể xóa section này.

% [Điền mô tả Dashboard vào đây]

\begin{figure}[H]
  \centering
  \framebox[14cm][c]{\rule{0pt}{6cm} [Chèn screenshot Dashboard]}
  \caption{Giao diện Dashboard báo cáo}\label{fig:dashboard}
\end{figure}
